% ---------------------------------------------------------------------
\section{How It Works (Architecture)}
\label{sec:arch}
% ---------------------------------------------------------------------

\subsection{The normal path}
\begin{enumerate}[leftmargin=1.6em,nosep]
  \item The user sends a prompt to the \textbf{Go controller} over HTTP.
  \item The controller knows which worker holds which layers and sends the
  prompt down the pipeline.
  \item Worker~1 runs its layers, passes activations to worker~2, and so on.
  \item The last worker produces tokens, which stream back to the user.
  \item Throughout, every worker sends a \textbf{heartbeat} to the controller.
\end{enumerate}

\subsection{The failover path (the novel core)}
\begin{enumerate}[leftmargin=1.6em,nosep]
  \item A worker stops sending heartbeats (lid closed / Wi-Fi dropped / killed).
  \item The controller detects the gap within a short timeout.
  \item It picks a healthy machine that holds (or can load) the dead worker's
  layer slice --- preferring the identical hot standby.
  \item It reassigns the slice and \textbf{re-runs the current request}.
  \item Token streaming resumes; the user sees a short pause, not a crash.
\end{enumerate}

\subsection{Text diagram}
\begin{lstlisting}[style=plain]
                 +-------------------+
   user  ---->   |  Go CONTROLLER    |   heartbeats <---- all workers
   prompt        |  - task queue     |   failover logic lives here
                 |  - who-has-what   |
                 +---------+---------+
                           | assign layer slices
         +-----------------+-----------------+
         v                 v                 v
   +-----------+     +-----------+     +-----------+
   | WORKER 1  | --> | WORKER 2  | --> | WORKER 3  | --> tokens
   | 4050 (win)|     | 4050 (win)|     | 3050 (lin)|
   | layers0-8 |     | layers9-16|     |layers17-24|
   +-----------+     +-----------+     +-----------+
                      (hot standby
                       for worker 1)

   If WORKER 1 dies: controller detects silence ->
   reassigns layers 0-8 to the standby -> re-runs request ->
   tokens keep streaming.
\end{lstlisting}

\subsection{Why the controller, not llama.cpp, owns recovery}
llama.cpp's RPC mode aborts the process on any peer failure (\VER,
\S\ref{sec:landscape}). By keeping llama.cpp \emph{local} on each worker and
putting all coordination in the Go controller, a dead worker becomes a
controller event to handle --- not a fatal crash of the whole system.
